% LIVRABLE — document professionnel, traçable et proportionné au besoin.
% Copier dans le dossier du devoir ; supprimer les sections optionnelles inutiles.
\documentclass[11pt,a4paper]{article}

\newcommand{\chargerpreambule}[1]{%
  \IfFileExists{../../templates/#1}{\input{../../templates/#1}}{%
  \IfFileExists{../templates/#1}{\input{../templates/#1}}{%
  \IfFileExists{templates/#1}{\input{templates/#1}}{\input{#1}}}}}
\chargerpreambule{preamble-commun.tex}
\chargerpreambule{preamble-livrable.tex}
\usepackage{xurl}
\contextedevoir{ACX401.1 — ECNSA}{Chapitre [numéro] — [intitulé]}

\newcommand{\versiondocument}{1.0}
\newcommand{\organisationetudiee}{[Organisation ou laboratoire]}

\begin{document}

\fairepagegarde{Livrable de sécurité}{[Titre du livrable]}%
  {\organisationetudiee — [objet précis du document]}%
  {[Auteur(s) — classe]}{\today}

\tableofcontents
\newpage

\section{Synthèse exécutive}

\begin{blocpedagogique}{À lire en premier}
\textbf{Objet :} [Quel besoin ce document traite-t-il ?]\\
\textbf{Résultat principal :} [Que montre l’analyse ?]\\
\textbf{Priorité :} [Quelle action faut-il décider ou mener d’abord ?]
\end{blocpedagogique}

Écrire cette synthèse pour une personne qui ne lira peut-être que cette page.
Développer les preuves, les choix et les limites dans les sections suivantes.

\filconducteur

\section{Objet, périmètre et méthode}

\subsection{Contexte et objectif}

[Présenter le besoin, l’objectif du document et le résultat attendu. Expliquer
les termes importants à leur première utilisation.]

\subsection{Périmètre}

\tableauECNSA
\begin{tabularx}{\linewidth}{>{\bfseries}p{3.2cm} X}
  \toprule
  Élément & Description \\
  \midrule
  Inclus & [Systèmes, services, personnes, sites, dates ou données couverts.] \\
  Hors périmètre & [Éléments exclus et justification.] \\
  Cadre / autorisation & [Environnement de test et limites autorisées.] \\
  Sources & [Documents, captures, commandes ou entretiens utilisés.] \\
  \bottomrule
\end{tabularx}
\fintableauECNSA

\subsection{Démarche}

[Décrire les étapes dans l’ordre : collecte, analyse, vérification, classement.
Indiquer les outils utiles et les choix méthodologiques.]

\section{Résultats en bref}

\small
\tableauECNSA
\begin{tabularx}{\linewidth}{>{\bfseries}p{2.1cm} X X >{\raggedright\arraybackslash}p{2.3cm}}
  \toprule
  ID / priorité & Constat vérifié & Conséquence possible & Suite proposée \\
  \midrule
  C-01 / élevée & [Fait observé et vérifiable.] & [Impact formulé simplement.] & [Action prioritaire.] \\
  C-02 / modérée & [Fait observé.] & [Impact possible.] & [Action ou vérification.] \\
  \bottomrule
\end{tabularx}
\fintableauECNSA
\normalsize

\section{Constats détaillés et preuves}

% Réutiliser ce bloc pour chaque constat important. Séparer observation et analyse.
\subsection{C-01 — [Titre factuel du constat]}

\textbf{Observation.} [Décrire le fait mesuré, vu ou trouvé. Éviter les
interprétations dans cette phrase.]

\textbf{Preuve.} [Référence précise à une capture, un tableau, un fichier,
une commande ou une source. Vérifier que la preuve est lisible et expurgée.]

\textbf{Explication.} [Dire en termes simples ce que montre la preuve, pourquoi
c’est important et quel actif ou besoin est concerné.]

\textbf{Risque ou impact.} [Décrire un scénario plausible, ses conséquences
et, si utile, les conditions nécessaires. Ne pas présenter une hypothèse
comme un fait.]

\textbf{Limite.} [Indiquer ce que le travail ne permet pas de confirmer.]

\captureanalyse{fig:preuve-1}{images/preuve-1.png}%
  {[Filtre Wireshark, commande ou méthode exacte]}%
  {[La légende dit précisément ce que la preuve établit.]}

\section{Recommandations et plan d’action}

Une bonne recommandation répond à un constat identifié, nomme une action
réalisable et indique comment vérifier qu’elle a fonctionné.

\small
\tableauECNSA
\begin{tabularx}{\linewidth}{>{\bfseries}p{1.4cm} X X >{\raggedright\arraybackslash}p{2.1cm} >{\raggedright\arraybackslash}p{1.8cm}}
  \toprule
  Priorité & Action et résultat attendu & Vérification / preuve & Responsable & Échéance \\
  \midrule
  P1 & [Action liée à C-01.] & [Contrôle observable.] & [Rôle.] & [Date / délai.] \\
  P2 & [Action liée à C-02.] & [Contrôle observable.] & [Rôle.] & [Date / délai.] \\
  \bottomrule
\end{tabularx}
\fintableauECNSA
\normalsize

\section{Limites et suites}

[Documenter les données manquantes, les vérifications non réalisées et les
prochaines étapes qui permettraient d’améliorer ou de confirmer les résultats.]

\section{Conclusion}

[Rappeler le résultat, les priorités et les décisions attendues. La conclusion
ne doit pas introduire de nouvelle preuve.]

\section*{Références}
\addcontentsline{toc}{section}{Références}

\begin{itemize}[leftmargin=*]
  \item [Auteur ou organisme, titre, version ou date, page ou adresse.] 
\end{itemize}

\appendix
\section{Annexes (facultatif)}

[Ajouter uniquement les éléments utiles à la vérification : extraits de
configuration expurgés, tableaux détaillés, glossaire ou preuves complémentaires.]

\end{document}
